% Drop-in for the experiments/setup of the codebook-design study. Explains and
% justifies the block-32 configuration vs. the deployed block-16 + full-Hadamard
% pipeline.
\paragraph{Block size.} The deployed GF4 pipeline (Section~\ref{sec:sw-support})
quantizes with a per-block scale over $16$ elements and rotates each layer by a
full $\mathrm{next\_pow2}(d_{\mathrm{in}})$ randomized Hadamard. For the
codebook-design study we instead fix a single block of $32$ with a block-diagonal
Hadamard. This is deliberate. The codebook is optimized on the RMS-\emph{normalized}
magnitude $m=\mathrm{clip}(|z|/\alpha,0,1)$ (Eq.~\ref{eq:codebook-opt}), where the
per-block scale absorbs the dynamic range and only the \emph{shape} of the
distribution enters the objective; that shape is approximately half-normal after
any randomized Hadamard, essentially independent of the block width, so the solved
levels are insensitive to the block granularity. Fixing block $32$ therefore
removes block size as a confound while giving one controlled configuration across
the eight-model, three-family suite, and keeps the per-model calibrate-and-solve
inexpensive. The universality of the result is direct evidence of this
insensitivity: the codebook solved independently on each model's own activations
is the same $k/16$ lattice across all layers of a model and across models and
families (Table~\ref{tab:codebook-suite}, Figure~\ref{fig:codebook-suite}), so the
optimum does not track the block choice. We consequently report the study at
block $32$ and deploy the identical codebook, unchanged, in the block-$16$
pipeline.
